% ---
% id: meat-core/rem:basis-invariance
% kind: rem
% title: "Invariance under the declared basis of the Choi input"
% summary: >-
%   Within one support frame the declared isometry is fixed up to a unitary on
%   A₀, which carries the Choi program onto itself, so the reported estimate,
%   h_lb and key length are invariant. Across frames only h_low and the sector
%   report are frame-free, and h_lb ≤ h_low holds in each frame.
% prominence: standard
% status: accepted
% uses: [lem:choi-marginal, def:support-frame, thm:choi, def:feasible, def:homog, rem:ordering, prop:frame-free-rows]
% mentions: [meat-core/frank-wolfe/thm:firstorder, def:branch-support, cor:full-gradient, req:frame-recorded]
% assumes: []
% refs: []
% provenance: newly-derived
% lean:
% ---

\begin{remark}
\label{rem:basis-invariance}
In trusted-source mode the isometry $\iota$ of \eqref{eq:iota} is not
canonical --- neither its range, the support frame, nor, within a fixed
range, its columns --- and this module declares no canonical choice. It
needs none, and the two freedoms are settled separately. \emph{Within one
range} $\iota$ is canonical up to a unitary on $A_0$, and every quantity
this module computes is invariant under that freedom; that is (a)--(d).
\emph{Across ranges} the programs genuinely differ --- they have
different numbers of variables --- but their value does not: by
Theorem~\ref{thm:choi} each maps \emph{onto} the one set $\Feas$, whose
infimum is $\hlow$ by \eqref{eq:hlow-def}, a definition that never
mentions $A_0$. What is frame-free across ranges is therefore the
\emph{mathematical} $\hlow$ of \eqref{eq:hlow-def}, and not the reported
numbers: the reported \tref{estimate} is the objective at the returned
\emph{point}, and the reported key length is computed from
$\max\{0,\hlb\}$, so both are functions of the iterates and neither
\emph{is} $\hlow$. Across ranges the programs carry different numbers of
variables, a run in one frame is not the image of a run in the other, and
those three reported numbers --- the estimate, $\hlb$ and the key
length --- may all differ. What (a)--(d) establish is the
\emph{within-range} statement, and there they are invariant, the two runs
corresponding under $\mathrm{Ad}_{U}$. What holds in every frame
regardless is the inequality: $\hlb\le\hlow$ by
Theorem~\ref{thm:firstorder}\,(ii), respectively~(iiib), applied in that
frame, against the one frame-free $\hlow$.

\emph{(a) Within one range the freedom is exactly a unitary.} Let
$\iota,\iota'$ both satisfy \eqref{eq:iota} \emph{with the same range},
write $\Pi\coloneqq\iota\iota^{\dagger}=\iota'\iota'^{\dagger}$ and put
$U\coloneqq\iota^{\dagger}\iota'$. Since
$\operatorname{ran}\iota'=\operatorname{ran}\Pi$ we have
$\Pi\iota'=\iota'$, so
$U^{\dagger}U=\iota'^{\dagger}\iota\,\iota^{\dagger}\iota'
=\iota'^{\dagger}\Pi\iota'=\iota'^{\dagger}\iota'=\idon{A_0}$, and
likewise $UU^{\dagger}=\iota^{\dagger}\Pi\iota=\idon{A_0}$: $U$ is unitary
on $A_0$, and
$\iota U=\iota\,\iota^{\dagger}\iota'=\Pi\iota'=\iota'$. Conversely
$\iota U$ satisfies \eqref{eq:iota}, with the same range, for every
unitary $U$ on $A_0$. The two admissible Choi factors are therefore
related by $W'_{\tau}\coloneqq\tau_A^{1/2}\iota'=W_{\tau}U$, and every
identity of Lemma~\ref{lem:choi-marginal}\,(i)--(iii) holds verbatim for
either:
$W'_{\tau}W'^{\dagger}_{\tau}
=W_{\tau}UU^{\dagger}W_{\tau}^{\dagger}=\tau_A$, and
$W'^{+}_{\tau}=(W_{\tau}U)^{+}=U^{\dagger}W^{+}_{\tau}$ (the
Moore--Penrose inverse of a product with a unitary right factor), so
$W'_{\tau}W'^{+}_{\tau}=W_{\tau}W^{+}_{\tau}
=\Pi_{\operatorname{supp}\tau_A}$ and
$W'^{+}_{\tau}W'_{\tau}=U^{\dagger}P_{\iota}U$, which is $\idon{A_0}$
exactly when the frame is tight. In particular $\ker W'_{\tau}
=U^{\dagger}(\ker W_{\tau})$ has the same dimension, so tightness is a
property of the range and not of the columns.

\emph{(b) The two parametrizations differ by a relabelling.} For a unitary
$U$ on $A_0$ write
$\mathrm{Ad}_{U}(J)\coloneqq(U\otimes\idon{})\,J\,
(U^{\dagger}\otimes\idon{})$ on $\Lin{A_0B'}$. From
$W'_{\tau}=W_{\tau}U$,
\begin{equation}\label{eq:chi-relabel}
  \chi'_{B'}=\chi_{B'}\circ\mathrm{Ad}_{U}.
\end{equation}
Moreover $\mathrm{Ad}_{U}$ carries
$\{J\in\Pos{A_0B'}:\tr_{B'}J=\idon{A_0}\}$ onto itself, with inverse
$\mathrm{Ad}_{U^{\dagger}}$: it preserves positivity, and
$\tr_{B'}\mathrm{Ad}_{U}(J)=U(\tr_{B'}J)U^{\dagger}
=U\idon{A_0}U^{\dagger}=\idon{A_0}$. Hence
Lemma~\ref{lem:choi-marginal} holds for either choice \emph{with the same
image} $\{\omega\in\Pos{AB'}:\tr_{B'}\omega=\tau_A\}$, and the two
parametrizations of one and the same $\omega$ are related by
$J'=\mathrm{Ad}_{U^{\dagger}}(J)$.

\emph{(c) The Choi program is carried onto itself.} Apply
$\mathrm{Ad}_{U}$ blockwise to $(J_s)_{s}$. Positivity and the coupling
constraint $\sum_s\tr_{\Qsys_s}J_s=\idon{A_0}$ are preserved by (b). For a
constraint entry, $(\idon{A_0}\otimes\hat O^{(s)}_k)$ commutes with
$U\otimes\idon{}$, so
\begin{equation*}
  \sum_s\tr_{\Qsys_s}\big[(\idon{A_0}\otimes\hat O^{(s)}_k)\,
  \mathrm{Ad}_{U}(J_s)\big]
  =U\Big(\sum_s\tr_{\Qsys_s}\big[(\idon{A_0}\otimes\hat O^{(s)}_k)\,
  J_s\big]\Big)U^{\dagger},
\end{equation*}
and conjugating by $W_{\tau}$ gives
$W_{\tau}U(\cdot)U^{\dagger}W^{\dagger}_{\tau}
=W'_{\tau}(\cdot)W'^{\dagger}_{\tau}$: the constraint of
\eqref{eq:choi-feas} built with $\iota'$ and imposed on $(J_s)_s$ is the
constraint built with $\iota$ and imposed on
$(\mathrm{Ad}_{U}(J_s))_s$. So $\mathrm{Ad}_{U}$ is a bijection from
$\Feas^{J}$ built with $\iota'$ onto $\Feas^{J}$ built with $\iota$. It
preserves the objective as well, because every ingredient of
\eqref{eq:choi-program} is a pullback along $\chi$: by
\eqref{eq:chi-relabel},
$P'^{J}_{\hat c}(J)=P^{J}_{\hat c}(\mathrm{Ad}_{U}(J))$,
$f'^{J}_{\hat c,s}(J_s)=f^{J}_{\hat c,s}(\mathrm{Ad}_{U}(J_s))$ and
$K'^{J}_{\hat c}(J)=K^{J}_{\hat c}(\mathrm{Ad}_{U}(J))$. The two programs
therefore have the same value, and $\hlow$ is unchanged. (Independently:
$\hlow$ is \emph{defined} by \eqref{eq:hlow-def} as an infimum over
$\Feas$ in the state variable $\omega$, which never mentions $A_0$;
Theorem~\ref{thm:choi} asserts that the Choi program equals that number
for any admissible $\iota$, so the invariance is already forced by the
theorem.)

\emph{(d) Everything downstream is invariant or covariant.} The
frame-free form \eqref{eq:frame-free-row} of a homogeneous row
(Proposition~\ref{prop:frame-free-rows}) is preserved, since
$c_k\idon{A_0}=U(c_k\idon{A_0})U^{\dagger}$ and conjugation by a unitary
preserves each of $\preceq$, $=$, $\succeq$. The branch operators
$X_{\hat c,s}=\cG_{\hat c,s}(\omega_s)$, the
\tref[maximal branch support]{maximal branch supports}
$W_{\hat c,s}$ of \eqref{eq:Wcs} and every value the objective takes are
\emph{literally} unchanged, being functions of
$\omega_s=\chi_{\Qsys_s}(J_s)$ and of $\operatorname{supp}\tau_A$ alone.
The gradients of Corollary~\ref{cor:full-gradient} are covariant:
$W'^{\dagger}_{\tau}M_xW'_{\tau}
=U^{\dagger}(W^{\dagger}_{\tau}M_xW_{\tau})U$ and
$\chi'^{\dagger}_{\Qsys_{s}}
=\mathrm{Ad}_{U^{\dagger}}\circ\chi^{\dagger}_{\Qsys_{s}}$, so both
$\Gamma^{(s')}_{\hat c}$ and $\Delta^{(s')}_{\hat c}$ transform as
$G\mapsto(U^{\dagger}\otimes\idon{})G(U\otimes\idon{})$, and a
Frank--Wolfe run in one declared basis is the
$\mathrm{Ad}_{U^{\dagger}}$-image of the same run in the other.
Consequently, \emph{within one range}, the reported estimate, $\hlb$ and
the reported key length are independent of which columns were chosen: the
two runs correspond under $\mathrm{Ad}_{U}$ and the objective, the
gradients and the pairings are carried along with them. Only the reported
$J_s$ and the channel reading of Remark~\ref{rem:ordering} are
basis-relative, and they are what the frame is recorded for
(Requirement~\ref{req:frame-recorded}). The sector report is frame-free in
the stronger sense, across ranges as well: its predicates read the declared
blocks $F^{(s)}_v$, $\hat O^{(s)}_k$ and $\{M_x\}$ and the declared tables,
and no $A_0$ enters any of them. \emph{Across ranges the reported numbers
are not claimed invariant}; what survives there is the mathematical
$\hlow$, and $\hlb\le\hlow$ in each frame, as the opening paragraph says.
\end{remark}
